\chapter{The simply-typed lamdba calculus}

So we don't want to allow self-reference. What do we do?
Russell enriched set theory with a type system to escape from the paradox, only allowing well-founded sets. Church thinks something similar could eliminate paradoxes from the
lambda-calculus, forbidding self-application in terms like $\pilambda x. xx$, and keeping the intuition
that lambda-abstractions denote functions.

If everything is a function, the problem is that we are allowing functions to take as parameters objects of any kind, for example another function.

The idea behind type theory is to classify lambda terms according to how they are meant to behave. We call these classifications \emph{types}, and we say that a term has a certain type and denote it as $t : T$ where $t$ is a term and $T$ is a type.

How can we decide the types of terms? Well remember that what we want to do is write mathematics. Recall the lambda terms that we already constructed in the previous chapter, and what we meant for them to represent

For example:

TRUE and FALSE represent boolean values,
the Church numerals 0,1,2,… represent natural numbers,
and terms such as EVEN represent computations on these objects.

Intuitively, these objects should not all behave in the same way. A boolean, a natural number, and a function computing parity play very different computational roles.

In particular, if EVEN is a function that determines whether a natural number is even, then we expect it to:

take a natural number as input,
and produce a boolean as output.

Thus, we would like to distinguish between terms such as:

TRUE,
0,
and EVEN

by assigning them different types.

We would like terms representing naturals to have a type representing natural numbers, for example type Nat:
$$
0 : \text{Nat}, \quad 1 : \text{Nat}, \quad \dots
$$
and true and false to have a type representing booleans, we can call this Bool.

Then, naturally the type of EVEN would be "something that takes a Nat and returns a Bool". This is written as
$$
EVEN : Nat \to Bool
$$

Of course, our definition of booleans and naturals were also functions, and thus they should have a type of the form $A \to B$ as well (and they do!) but let's not worry about this just now (as we don't yet have the language to precisely describe the types of these terms). For now, consider the following formal definition of types:

Definition (simple type): (simple) types are defined inductively as terms were; we have an infinite set of type variables, which we will denote (for example) $\alpha, \beta, \gamma, \dots$. Then
\begin{itemize}
  \item Base case: if $\alpha$ is a type variable then it is a type
  \item Recursive case: if $\alpha$ and $\beta$ are types then $\alpha \to \beta$ is also a type, called the arrow type.
\end{itemize}

Examples of simple types: $\alpha$, $\tau_1$, $\tau_1 \to \tau_2$, $(\alpha \to \beta) \to \gamma$, $\alpha \to \beta \to \gamma$.

Note that, as we did with terms, types are associated automatically to the right; so
$$
\alpha \to \beta \to \gamma = \alpha \to (\beta \to \gamma) \neq (\alpha \to \beta) \to \gamma
$$

Note that, up until now, the definition of types is independent from the lambda calculus. The notion of types was actually introduced by Russell in 1908 (so even before Church's lambda calculus) as an idea to solve his own paradox. The idea was that different level of objects could not be treated the same (a set of natural numbers must be different from a set of sets of natural numbers).

However the exact idea of Russell is not exactly this as this is more computational and already with the idea of putting it on top of the lambda calculus.

But I digress.

So now we want to put the types on top of our lambda terms. In particular, we define type assingments as:

Definition: given a lambda-term $t$ and a simple type $\tau$, a type assigment is an expression $t : \tau$.

Moreover, a type context is a finite set of assingments in which the term $t$ is in fact a variable, i.e. a set
$$
\Gamma = \{x_1 : \tau_1, \dots, x_n : \tau_n\},
$$
and given a context $\Gamma$, a term $t$ and a type $\tau$, a type judgement is an expression $$\Gamma \vdash t : \tau,$$ which is read as "under the assumptions $\Gamma$, the term $t$ can be assigned the type $\tau$".


Now even though this definition explains the valid syntax for type judgements, not all type judgements are true / vlaid themselves.

Valid judgements must be obtain by deriving from the following rules (which can be seen in a way as an inductive definition):

\begin{itemize}
  \item An infinite set of axioms: if $x$ is a (term) variable, then
  $$
  x : \tau \vdash x : \tau
  $$
  \item Deduction rules:
  \begin{itemize}
    \item (Abstraction) If assuming $\Gamma$ and $x : \tau$,we can derive that $t : \tau'$, then just assuming $\Gamma$ we can derive $\lambda x. t : \tau \to \tau'$
    \begin{prooftree}
      \AxiomC{$\Gamma \cup x : \tau \vdash t : \tau'$}
      \RightLabel{[ABS]}
      \UnaryInfC{$\Gamma \vdash \lambda x. t : \tau \to \tau'$}
    \end{prooftree}
    \item (Application) yada yada
    \begin{prooftree}
      \AxiomC{$\Gamma \vdash t : \tau \to \tau'$}
      \AxiomC{$\Gamma' \vdash t' : \tau$}
      \RightLabel{[APP]}
      \BinaryInfC{$\Gamma \cup \Gamma' \vdash tt' : \tau'$}
    \end{prooftree}
  \end{itemize}  
\end{itemize}

Examplessssssssssss

---

Type judgements. What are they. Why do we care about doing type judgements? How does this relate to being a "foundation of mathematics"? Difference between CIC and ZFC: we dont care about membership but about type judgements. Type judgements are described via rules (not axioms). What are these rules? Where do these rules come from? What problem do they solve?



What can we do with types? What is interesting about this for Lean/computation? In which way do we prefer types over sets? Type judgements over membership?

Are they enough? Why not? Why are we missing? What can we do with ZFC that we cannot do with simply typed lambda calculus? How do we solve this lacking?

\section{Type Judgements in Lean}

the check command? should i talk here about the elaborator??

Essentially the elaborator is the compiler of Lean. 

Here we  would write a 

\section{Exercises}

